\documentclass[12pt]{article}
\usepackage{graphicx,amsmath,amsfonts,amssymb,epsfig,euscript,enumerate}
\usepackage[T1]{fontenc}
\usepackage[utf8x]{inputenc}

\newtheorem{exercise}{Ejercicio}
\newcommand{\bej}{\begin{exercise}\rm}
\newcommand{\fej}{\end{exercise}}

\newcommand{\R}{\mathbb{R}}
\newcommand{\C}{\mathbb{C}}
\def\dt{\Delta t}
\def\dx{\Delta x}

\topmargin-2cm \vsize 29.5cm \hsize 21cm
\setlength{\textwidth}{16.75cm}\setlength{\textheight}{23.5cm}
\setlength{\oddsidemargin}{0.0cm}
\setlength{\evensidemargin}{0.0cm}

\begin{document}
\centerline{{\small Universidad de Buenos Aires - Facultad de Ciencias Exactas y Naturales - Depto. de Matemática}}
 
 \vskip 0.2cm
 \hrulefill
 \vskip 0.2cm

 \centerline{{\bf\Huge {\sc Elementos de Cálculo Numérico}}}
 \vskip 0.2cm
 \centerline{\ttfamily Primer Cuatrimestre 2026}
 \hrulefill

 \bigskip
 \centerline{\bf Ejercicios tipo parcial -- Guías 3}
 \bigskip

\begin{exercise}[Interpolación en distintas bases]
Dado un conjunto de datos $(x_i,y_i), i=0,\dots,n$, se propone hallar el polinomio interpolador de grado a lo sumo $n$. Para ello, se consideran las siguientes bases de polinomios:
\begin{itemize}
\item Base monomial: $\{1,x,x^2,\dots,x^n\}$.
\item Base de Lagrange: $\{L_0(x),L_1(x),\dots,L_n(x)\}$, donde $L_i(x_j)=\delta_{ij}$.
\item Base de Newton: $\{1,(x-x_0),(x-x_0)(x-x_1),\dots,(x-x_0)(x-x_1)\cdots(x-x_{n-1})\}$.
\end{itemize}


\begin{enumerate}[(a)]
\item ¿Cuál es el sistema lineal que se debe resolver para hallar el polinomio interpolador en cada una de las bases mencionadas?
\item En la base de Lagrange, ¿cómo pueden obtenerse los coeficientes del polinomio interpolador directamente a partir de los datos $(x_i,y_i)$, es decir, sin necesidad de resolver numéricamente un sistema lineal?
\item ¿Cuál es el costo computacional del procedimiento en cada caso?
\end{enumerate}
\end{exercise}


\begin{exercise}[Interpolación en distintos nodos]
Considere la función $f(x) = x^7$. Se desea aproximar $f$ mediante un polinomio de grado 3, en el intervalo $[-1,1]$. Para ello, se consideran dos conjuntos de nodos:
\begin{itemize}
\item Nodos equiespaciados: $x_i = \lbrace -1, -\frac{1}{3}, \frac{1}{3}, 1\rbrace$.
\item Nodos $x_i = \lbrace-1, -\frac{1}{\sqrt{5}}, \frac{1}{\sqrt{5}}, 1\rbrace$.
\end{itemize}
\begin{enumerate}[(a)]
\item Utilizando la base de Lagrange, encuentre el polinomio interpolador de grado 3 para cada conjunto de nodos.
\item Compare las aproximaciones obtenidas en el punto $x=0.1$. ¿Cuál es más precisa?
\end{enumerate}
\end{exercise}


\end{document}
